% LaTeXbangla user manual. Copyright 2016, 2026 Adib Hasan.
% Distributed under the LPPL, version 1.3c or later.
\documentclass[11pt,a4paper]{article}
\usepackage[margin=22mm,headheight=15pt]{geometry}
\usepackage{latexbangla}
\setmainfont{TeX Gyre Pagella}
\setsansfont{TeX Gyre Heros}
\setmonofont{TeX Gyre Cursor}
% Give the opening example a visibly different English family. Only its
% preview uses this family; the surrounding manual stays in Pagella.
\newfontfamily\manualexampleenglishfont{TeX Gyre Adventor}

% Save a configured family for the boldness preview; the main text uses the
% normal package default. Font setters belong in the preamble, including here.
\setbanglamainfont[AutoFakeBold=5,AutoFakeSlant=0.2]{Kalpurush}
\NewCommandCopy\manualstrongfont\bengalifont
% Configure the family comparison with the public setters, then save its
% families. The preview selects them locally; prose and listings keep Kalpurush.
\setbanglamainfont{Noto Sans Bengali}
\NewCommandCopy\manualsansfont\bengalifont
\setbanglamainfont{Siyam Rupali}
\NewCommandCopy\manualmonofont\bengalifont
\setbanglamainfont{Kalpurush}

\usepackage{amsmath,amsthm,graphicx,booktabs,tabularx}
\usepackage{xcolor,minted,tcolorbox,titlesec,fancyhdr,needspace}
\tcbuselibrary{raster,skins}
\usepackage{hyperref}
\usepackage{alphalph}
% The extended-label example uses alphalph's custom-alphabet interface.
\makeatletter
\newalphalph{\extendedBangla}{\@alph}{36}
\makeatother
\definecolor{BanglaGreen}{HTML}{006847}
\definecolor{BanglaInk}{HTML}{263640}
\definecolor{BanglaRed}{HTML}{C92935}
\definecolor{CodePaper}{HTML}{F4F6F8}
\definecolor{ResultPaper}{HTML}{F0F8F4}
\hypersetup{colorlinks,linkcolor=BanglaGreen,urlcolor=BanglaGreen,
  citecolor=BanglaGreen,pdftitle={LaTeXbangla 1.1: User guide},pdfauthor={Adib Hasan}}
\setlength{\parindent}{0pt}
\setlength{\parskip}{4pt plus 1pt}
\setlength{\emergencystretch}{2em}
\titleformat{\section}{\Large\sffamily\bfseries\color{BanglaGreen}}{\thesection}{0.6em}{}
\titleformat{\subsection}{\normalsize\sffamily\bfseries\color{BanglaInk}}{}{0pt}{}
\titlespacing*{\section}{0pt}{9pt}{6pt}
\titlespacing*{\subsection}{0pt}{8pt}{3pt}
% Reserve room before starting a heading so both its text and PDF anchor move
% together when a page is full. Sections otherwise flow without forced breaks.
\AddToHook{cmd/section/before}{\Needspace{5\baselineskip}}
\AddToHook{cmd/subsection/before}{\Needspace{6\baselineskip}}
\pagestyle{fancy}
\fancyhf{}
\fancyhead[L]{\small\sffamily\color{BanglaGreen}LaTeXbangla}
\fancyhead[R]{\small\sffamily\color{BanglaInk}Version 1.1\quad User guide}
\fancyfoot[C]{\small\thepage}
\renewcommand{\headrulewidth}{0.3pt}
% Keep the English guide's navigation Western while examples use the package's
% default Bengali numbering. Raw counter values bypass localized formatters.
\renewcommand{\thesection}{\number\value{section}}
\renewcommand{\thesubsection}{\thesection.\number\value{subsection}}
\renewcommand{\thepage}{\number\value{page}}
\newcommand{\pkg}[1]{\texttt{#1}}
\newtheorem{theorem}{উপপাদ্য}
\providecommand{\theoremautorefname}{উপপাদ্য}
\renewcommand{\theoremautorefname}{উপপাদ্য}

\tcbset{manualbox/.style={enhanced,boxrule=0.4pt,arc=2mm,
  colframe=BanglaGreen!35!white,colback=CodePaper,
  colbacktitle=BanglaGreen!9!white,coltitle=BanglaGreen,
  fonttitle=\small\sffamily\bfseries,left=3mm,right=3mm,top=2mm,bottom=2mm,
  before skip=5pt,after skip=5pt}}
\tcbset{manualwarning/.style={manualbox,
  colframe=BanglaRed!45!white,colback=BanglaRed!5!white}}
% Minted supplies LaTeX syntax highlighting through Pygments. Its verbatim
% renderer preserves the Bangla text and uses the same fonts as the document.
\setminted{style=default,breaklines=true,breakanywhere=true}
\newminted[code]{latex}{fontsize=\small,samepage}

% Write each example once, then show its source and execute that same file.
% A plain integer in the filename keeps generated names independent of the
% document's numbering policy. Each output panel is scoped to its own box.
\newcounter{manualexample}
\newcommand{\manualcodetitle}{Code}
\newcommand{\manualresulttitle}{Result}
\newcommand{\manualexamplefile}{\jobname-example-\number\value{manualexample}.tex}
\newenvironment{examplepair}
  {\VerbatimEnvironment\stepcounter{manualexample}%
   \begin{VerbatimOut}{\manualexamplefile}}
  {\end{VerbatimOut}%
   \begin{tcbraster}[raster columns=2,raster column skip=3mm,
     raster equal height=rows,raster valign=top]
   \begin{tcolorbox}[manualbox,title=\manualcodetitle]
     \inputminted[fontsize=\footnotesize]{latex}{\manualexamplefile}
   \end{tcolorbox}
   \begin{tcolorbox}[manualbox,title=\manualresulttitle,colback=ResultPaper]
     \setlength{\parskip}{4pt}%
     % Match normal document footnote numbering inside the boxed minipage.
     \renewcommand{\thempfootnote}{\arabic{mpfootnote}}%
     \input{\manualexamplefile}
   \end{tcolorbox}
   \end{tcbraster}}

\begin{document}
\thispagestyle{empty}
\begin{center}
\includegraphics[width=0.66\linewidth]{logo.png}\\[2pt]
{\Large\sffamily\color{BanglaGreen}Writing Bangla with ordinary \LaTeX}\\[5pt]
{Adib Hasan}\\[2pt]
{\small\href{https://groups.google.com/g/latexbangla}{\texttt{latexbangla@googlegroups.com}}}\\[3pt]
{23 October 2026\quad\textcolor{BanglaRed}{\textbullet}\quad Version 1.1}
\end{center}

LaTeXbangla is a package for writing Bangla articles, reports and books with
ordinary \LaTeX\ commands. To begin, load the package and choose a Bangla font.
You can then write Bangla and English together, make text bold with
\verb|\textbf|, and create lists, equations and references as usual. The package
selects the appropriate font for each script while leaving your document layout,
English fonts and mathematics fonts to you.

LaTeXbangla also takes care of the
surrounding document. Lists use labels such as ১ and (ক), automatic references
use names such as সমীকরণ, and numeric citations appear in Bangla. It builds on
fontspec and Polyglossia, adding automatic font switching, bold and italic
fallbacks, alphabetic numbering, Bangla math digits and reference support.
\hyperref[manual:bangla]{A complete Bangla translation follows this English guide}.

\section{Quickstart}
Compile with \textbf{XeLaTeX or LuaLaTeX} using TeX Live 2026 or a compatible newer release.
Install a Bangla font from the list in section~2. Add this preamble:
\begin{tcolorbox}[manualbox,title=Minimal document]
\begin{code}
\documentclass{article}
\usepackage{fontspec}
\setmainfont{TeX Gyre Adventor} % Optional; English font is preserved.
\usepackage{latexbangla}
\setbanglamainfont{Kalpurush}
\begin{document}
  % Write your document here.
\end{document}
\end{code}
\end{tcolorbox}
Put the following between \verb|\begin{document}| and \verb|\end{document}|.
Notice how the English keeps Adventor's geometric sans serif shapes, even inside
bold and italic text. Each Code/Result pair executes the source printed beside it.
\begingroup
\let\rmfamilylatin\manualexampleenglishfont
% Reselect after replacing the companion so both engines capture this local setup.
\bengalifont
\begin{examplepair}
বাংলা and English একই অনুচ্ছেদে।

English keeps its own typeface.

\textbf{গুরুত্বপূর্ণ: read this.}
\textit{বাংলা ও English একসঙ্গে।}
\end{examplepair}
\endgroup

\subsection{The same paragraph with fontspec alone}
For occasional Bangla in an English document, fontspec alone is often enough.
Here is a complete document that reproduces the paragraph above. Declare a
Bangla font, give Kalpurush its synthetic styles, and mark each Bangla passage
with that font. English continues to use Adventor outside those groups.
Under LuaLaTeX, also enable HarfBuzz for Bangla shaping:
\begin{tcolorbox}[manualbox,title=Without LaTeXbangla]
\begin{code}
\documentclass{article}
\usepackage{fontspec,iftex}
\setmainfont{TeX Gyre Adventor}
\ifluatex\defaultfontfeatures{Renderer=HarfBuzz}\fi
\newfontfamily\banglatext[
  Script=Bengali,
  AutoFakeBold=4,
  AutoFakeSlant=0.2
]{Kalpurush}
\begin{document}
{\banglatext বাংলা} and English
{\banglatext একই অনুচ্ছেদে।}

English keeps its own typeface.

\textbf{{\banglatext গুরুত্বপূর্ণ:} read this.}
\textit{{\banglatext বাংলা ও} English
  {\banglatext একসঙ্গে।}}
\end{document}
\end{code}
\end{tcolorbox}
LaTeXbangla supplies these font switches and style defaults, so you can write
the mixed paragraph directly. The example above matches the text's appearance;
the Bangla numbering and translated document labels described later in this
guide are additional parts of LaTeXbangla's language support.

\textbf{Install locally.} Copy \pkg{latexbangla.sty},
\pkg{latexbangla-bengali.lbx} and the complete \pkg{code/} folder beside your
main document, or install that structure in your TeX tree. Package files are in
\href{https://github.com/Neehan/LaTeXbangla}{the repository's \pkg{latexbangla/} folder}.
The tested versions are LaTeX 2026-06-01, Polyglossia 2.13 and fontspec 2.9h.

\section{Choose fonts and text styles}
\subsection{Supported Bangla fonts}
Choose a font that suits your document. Any family below can be your only
Bangla font; put its name in \verb|\setbanglamainfont{...}|, or use Anek's
filename as noted below. You can compare
samples in the font guide included in
\href{https://github.com/Neehan/LaTeXbangla}{the repository}.
\begin{center}
\begin{tabularx}{\linewidth}{@{}lX@{}}
\toprule
\textbf{Font name} & \textbf{Default Bangla bold} \\
\midrule
\texttt{Kalpurush} & Synthetic, strength 4 \\
\texttt{SolaimanLipi} & Synthetic, strength 4 \\
\texttt{Siyam Rupali} & Synthetic, strength 4.5 \\
\texttt{Noto Serif Bengali} & Real bold face \\
\texttt{Noto Sans Bengali} & Real bold face \\
\texttt{Anek Bangla} & Variable file: synthetic, strength 4 \\
\texttt{Tiro Bangla} & Synthetic, strength 3.5; real italic face \\
\texttt{Galada} & Synthetic, strength 3.5 \\
\texttt{FreeSerif} & Synthetic, strength 3.5; real italic face \\
\bottomrule
\end{tabularx}
\end{center}
LaTeXbangla uses bold and italic faces that contain Bangla glyphs, creating
missing styles with added weight or slant. For unlisted fonts, these synthetic
fallbacks use a weight of 2.5 and a slant of 0.2.
Other shapes, including small caps, depend on the font. Install fonts separately.

\begin{tcolorbox}[manualwarning]
If Anek Bangla fails to load by name, save its TTF as \texttt{AnekBangla-VF.ttf}
beside your document and use \verb|\setbanglamainfont{AnekBangla-VF.ttf}|.
\end{tcolorbox}

\subsection{Separate Bangla and English families}
Give body text, headings and typewriter text their own appearance while keeping
Bangla and English together. Here we pair Kalpurush with Pagella for the main
text, Noto Sans Bengali with Heros for sans serif, and Siyam Rupali with Cursor
for the mono role. Put these choices in the preamble, before \verb|\begin{document}|:
\begin{tcolorbox}[manualbox,title=Optional font choices]
\begin{code}
\setbanglamainfont{Kalpurush}
\setbanglasansfont{Noto Sans Bengali}
\setbanglamonofont{Siyam Rupali}
\setmainfont{TeX Gyre Pagella}
\setsansfont{TeX Gyre Heros}
\setmonofont{TeX Gyre Cursor}
\end{code}
\end{tcolorbox}
Now use \verb|\textrm|, \verb|\textsf| and \verb|\texttt| as usual. Each
selects the matching pair of fonts, including their bold and italic styles:
\begingroup
\let\bengalifontsf\manualsansfont
\let\bengalifonttt\manualmonofont
\begin{examplepair}
বাংলা ও English।\\
\textbf{গাঢ় ও bold।}\\
\textit{তির্যক ও italic।}

\textsf{বাংলা ও English।}\\
\textsf{\textbf{গাঢ় ও bold।}}\\
\textsf{\textit{তির্যক ও italic।}}

\texttt{বাংলা ও English।}\\
\texttt{\textbf{গাঢ় ও bold।}}\\
\texttt{\textit{তির্যক ও italic।}}
\end{examplepair}
\endgroup
The same text takes on three different appearances without marking individual
Bangla or English passages. Siyam Rupali keeps its natural character widths;
assigning it to the mono role does not make it a fixed-width font.

These choices are optional: one Bangla main font can fill all three roles, and
your English fonts stay as they are until you change them. Bangla and English
setters can appear in either order; each replaces its own role.
Automatic font switching does not change the text's language or hyphenation
rules. For an English passage with English hyphenation, declare
\verb|\setotherlanguage{english}| in the preamble and use \verb|\textenglish{...}|.

\subsection{Fontspec features, files and spacing}
To match another typeface's size, adjust word spacing or use a font file from
your project, pass fontspec options to any Bangla setter. The options can go
before or after the font name:
\begin{code}
\setbanglamainfont[features]{font}
\setbanglamainfont{font}[features]
\end{code}
The defaults include \pkg{Script=Bengali} and, under LuaLaTeX,
\pkg{Renderer=HarfBuzz}.
Your choices override the defaults. For example, adjust size and spacing or load a file:
\begin{code}
\setbanglamainfont[Scale=1.05,WordSpace=1.1]{Kalpurush}
\setbanglamainfont[Path=fonts/]{fontkalpurush.ttf}
\end{code}
Here \pkg{WordSpace=1.1} adds 10\% to the font's native word spacing,
replacing the preset. These are the default size and spacing settings:
\begin{center}
\begin{tabular}{lcc}
\toprule
Font & Scale & WordSpace \\
\midrule
Kalpurush & 1 & 1 \\
SolaimanLipi & 1 & 1.2 \\
Siyam Rupali & 0.8 & 1.12 \\
Noto Serif Bengali & 1 & 1.3 \\
Noto Sans Bengali & 1 & 1.2 \\
Anek Bangla & 0.9 & 2.1 \\
Tiro Bangla & 1 & 1.25 \\
Galada & 0.9 & 1.67 \\
FreeSerif & 1 & 1.33 \\
\bottomrule
\end{tabular}
\end{center}
Presets recognize family names and standard filenames; renamed files use
generic defaults. To choose faces from separate
files, use \pkg{UprightFont}, \pkg{BoldFont}, \pkg{ItalicFont} and \pkg{BoldItalicFont}.

\section{Control boldness and emphasis}
Use \verb|\textbf| to draw attention to a definition and \verb|\textit| for
emphasis; both work with the default setup. If a font's synthetic bold needs
more contrast in your document, increase \pkg{AutoFakeBold}.
Include \pkg{AutoFakeSlant} to carry that weight into bold italic too.
Here is how the four styles look with a bold strength of 5:
\begin{tcolorbox}[manualbox,title=Custom synthetic styles (preamble)]
\begin{code}
\setbanglamainfont[
  AutoFakeBold=5,
  AutoFakeSlant=0.2
]{Kalpurush}
\end{code}
\end{tcolorbox}
% This preview uses precisely the preamble settings printed above.
\begingroup
\let\bengalifont\manualstrongfont
\bengalifont
\begin{examplepair}
সাধারণ লেখা।

\textbf{গাঢ় লেখা।}

\textit{তির্যক লেখা।}

\textbf{\textit{গাঢ় ও তির্যক।}}
\end{examplepair}
\endgroup
Explicit \pkg{AutoFakeBold} requests synthetic bold even if a real bold face
exists. Use it with a numeric value: fontspec's bare \pkg{AutoFakeBold} means
1.5, rather than LaTeXbangla's font-specific default.

You can also adjust bold and bold italic separately through fontspec's face features:
\begin{code}
\setbanglamainfont[
  BoldFeatures={FakeBold=5},
  BoldItalicFeatures={FakeBold=5,FakeSlant=0.2}
]{Kalpurush}
\end{code}
These options are especially useful with a family that needs synthetic faces.
Keep the ordinary setter when you prefer the tested defaults and real-face
discovery. You can also configure Polyglossia's \verb|\bengalifont|,
\verb|\bengalifontsf| and \verb|\bengalifonttt| directly, supplying their styles
yourself. Use \verb|\renewfontfamily| for an existing command, including the
sans and mono defaults that LaTeXbangla defines.

\subsection{Punctuation, bullets and footnotes}
Punctuation and footnotes fit into mixed text in the usual way. Typed dashes, curly
quotes and ellipses come from the English family, which also supplies standard
bullets and symbolic footnote marks missing from the Bangla font. Ordinary
footnotes use Bangla numbers; Roman or symbolic formats keep their structure.
\begin{examplepair}
“বাংলা”—English… আবার বাংলা।
\begin{itemize}
  \item একটি কথা।
  \item আরেকটি কথা।
\end{itemize}
একটি তথ্য\footnote{ছোট ব্যাখ্যা।}
\end{examplepair}

\section{Numbers, lists and your own counters}\label{sec:numbering}
Write an exercise with subparts using ordinary nested lists. LaTeXbangla supplies
the Bangla labels, so the same list commands work throughout your document.

\pkg{changecounternumbering=true} is the default: decimal counters become
Bengali, and \verb|\alph| and \verb|\Alph| use Bengali letters. Values, reset
rules and compound labels keep their structure; Roman and symbolic formats stay
as chosen. This English part keeps section, subsection and page numbers
Western; its examples use the default Bangla numbering. The four standard
\pkg{enumerate} levels become ১., (ক), i. and ক.; their punctuation and
Roman numerals stay as in ordinary \LaTeX.
\begin{examplepair}
\begin{enumerate}
  \item একটি আয়তক্ষেত্রের দৈর্ঘ্য ৪,
    প্রস্থ ৩ একক।
  \begin{enumerate}
    \item ক্ষেত্রফল কত?
    \item কর্ণের দৈর্ঘ্য কত?
  \end{enumerate}
\end{enumerate}
\end{examplepair}

\subsection{Custom counters need no registration}
For numbered exercises, create a counter as usual. Use \verb|\refstepcounter|
to increment it and allow \verb|\label| and \verb|\ref|;
\verb|\stepcounter| increments it without setting up a reference.
\begin{examplepair}
\newcounter{exercise}
\setcounter{exercise}{11}
\refstepcounter{exercise}
\label{ex:custom}
অনুশীলনী \theexercise।\par
\renewcommand{\theexercise}
  {\alph{exercise}}
\setcounter{exercise}{2}
অনুশীলনী \theexercise।
\end{examplepair}
Here, \texttt{[section]} resets the counter at each section;
\verb|\theproblem| adds the section to its printed number:
\begin{code}
\newcounter{problem}[section]
\renewcommand{\theproblem}{\thesection.\arabic{problem}}
\end{code}
You can also choose the digit script for just one counter:
\begin{code}
\renewcommand{\theexercise}{\tobangla{exercise}}       % Bangla
\renewcommand{\theexercise}{\number\value{exercise}} % Western
\end{code}
A formatter that bypasses \verb|\arabic|, \verb|\alph| and \verb|\Alph|
keeps its own output.

\subsection{Keep selected counters in English}
You can keep Bangla numbering throughout the document and make an exception
for one counter. Its references follow the same format:
\begin{examplepair}
\renewcommand{\theexercise}
  {\number\value{exercise}}
\refstepcounter{exercise}
\label{ex:western}
Exercise \theexercise.\par
See Exercise~\ref{ex:western}.
\end{examplepair}
For Latin letters, load \verb|\usepackage{alphalph}| in the preamble.
Its \verb|\alphalph| formatter gives a, b, \ldots, z, aa, \ldots;
\verb|\AlphAlph| gives uppercase letters. Here the change applies to just
one list, because the definition is inside the environment:
\begin{examplepair}
\begin{enumerate}
  \renewcommand{\theenumi}
    {\alphalph{\value{enumi}}}
  \item First.\label{item:latin}
  \item Second.
\end{enumerate}
See item~\ref{item:latin}.
\end{examplepair}
For 1, 2, \ldots\ in that list, use \verb|\number\value{enumi}| instead.
For a custom counter, put \verb|\alphalph{\value{exercise}}| in
\verb|\theexercise|. Later lists keep their Bangla labels. These changes
affect the numbers; translated names supplied by \verb|\autoref| stay unchanged.

\subsection{Extended alphabetic labels}
Ordinary alphabetic labels stop at ৎ, the 36th letter. If you need more,
\pkg{alphalph} can continue with কক, কখ and so on. With counter conversion
enabled, add this to your preamble after loading LaTeXbangla:
\begin{code}
\usepackage{alphalph}
\makeatletter
\newalphalph{\extendedBangla}{\@alph}{36}
\makeatother
\end{code}
The last argument gives the alphabet size. Assign the new formatter to the
counter you want to extend; other counters keep their usual behavior:
\begin{examplepair}
\newcounter{extendedexercise}
\renewcommand{\theextendedexercise}
  {\extendedBangla{
    \value{extendedexercise}}}
\setcounter{extendedexercise}{36}
\theextendedexercise,
\stepcounter{extendedexercise}
\theextendedexercise,
\stepcounter{extendedexercise}
\theextendedexercise
\end{examplepair}
Use \verb|\refstepcounter| instead of \verb|\stepcounter| when you also need
\verb|\label| and \verb|\ref|. The built-in \verb|\alphalph| and
\verb|\AlphAlph| commands still produce Latin letters; use the custom
\verb|\extendedBangla| command for Bangla labels.

\subsection{Choose automatic or explicit conversion}
\verb|\usepackage[changecounternumbering=false]{latexbangla}| preserves the
existing counter and numeric citation formats while retaining fonts and
translations. This is a package option; language switches do not toggle it.
For a code with leading zeroes, or a Bangla counter value even when automatic
conversion is off, use:
\begin{examplepair}
\numtobangla{0012}; \tobangla{exercise}
\end{examplepair}
Both helpers are expandable. \verb|\numtobangla| takes Western digits only and preserves leading zeroes;
signs and decimal points are not accepted. \verb|\tobangla| prints a counter's
integer value, including zero and negative values.

\section{Mathematics, tables and theorems}
In a Bangla lesson, you can explain a formula in Bangla, keep familiar symbols
such as $x$ and $v$, and write numbers in either script.
Write numbers in the script you want to see: \texttt{12} stays 12 and
\texttt{১২} stays ১২, in text and in math. Automatic numbering leaves your
formulas alone. In math, Bangla digits come from your main Bangla font, sized to fit
their position in the formula; the other symbols use your mathematics font.

With \pkg{amsmath}, use \verb|\text| to name quantities and units in Bangla.
For a 12-metre journey taking 3 seconds, write:
\begin{examplepair}
\begin{equation}\label{eq:speed}
  \text{বেগ}=
  \frac{১২}{৩}=৪\text{ মি./সে.}
\end{equation}
\[
  x_৩^{২}+\sqrt{৪}=12+১২
\]
\end{examplepair}
The equation number is Bangla too, ready to cite with \verb|\eqref| later.
You can use the same Bengali digits in fractions, roots, matrices and nested scripts.
One-digit arguments can be unbraced, as in \verb|x^২|; use braces for several
digits. Use \verb|\boldmath| for bold mathematics: surrounding \verb|\textbf|
alone does not change math weight. Commands for special math alphabets, such as
\verb|\mathbb|, do not select alternate Bengali digit designs.
Bold versions of the other symbols depend on your mathematics font setup.

\subsection{Choose the mathematics font independently}
To give your formulas a different appearance, choose a font with \pkg{unicode-math}:
\begin{code}
\usepackage{unicode-math}
\setmathfont{TeX Gyre Pagella Math}
\end{code}
For traditional mathematics, load \pkg{amsmath,newtxmath} \emph{before}
\pkg{latexbangla}. The language package leaves math-font selection to you.

\subsection{Tables and author-defined environments}
Put Bangla descriptions beside English terms and measurements in an ordinary
table. Each cell uses the appropriate font automatically:
\begin{examplepair}
\begin{tabular}{llc}
বাহু & Side & দৈর্ঘ্য \\
\hline
প্রথম & First & ৩ \\
দ্বিতীয় & Second & ৪
\end{tabular}
\end{examplepair}
When using \pkg{siunitx} to format and align numbers, supply Western digits as
input; use a text column when you want to type Bangla digits directly.
Define theorems with \pkg{amsthm}, for example
\verb|\newtheorem{theorem}{উপপাদ্য}| in the preamble. Polyglossia supplies the
Bangla proof heading; your class or theorem style controls the design.
\begin{examplepair}
\begin{theorem}
\label{thm:pythagoras}
সমকোণী ত্রিভুজের লম্ব দুই বাহুর
দৈর্ঘ্য ৩ ও ৪ হলে অতিভুজের দৈর্ঘ্য ৫।
\end{theorem}
\begin{proof}
$৩^{২}+৪^{২}=৫^{২}$।
\end{proof}
\end{examplepair}

\section{Labels, references and citations}
Refer to an equation, theorem or exercise by its label, and the printed number follows
any later changes to the document. Here the ordinary reference commands point
back to the examples you have just seen.
Your class's chapter, figure, table, contents, appendix and bibliography names,
along with dates, come from Polyglossia's Bengali translations. With
\pkg{hyperref}, LaTeXbangla also gives automatic references Bangla names.
Compile again after adding labels
or citations so their references can be resolved.
\begin{examplepair}
\autoref{eq:speed};
\eqref{eq:speed}।

\autoref{thm:pythagoras}।
অনুশীলনী \ref{ex:custom}।
\end{examplepair}
\verb|\ref| prints the label, \verb|\eqref| adds parentheses, and
\verb|\autoref| includes the translated name. Names include পরিচ্ছেদ,
উপপরিচ্ছেদ, সমীকরণ and পাদটীকা. To choose another term, redefine its public
hyperref name after loading the packages:
\begin{code}
\renewcommand{\sectionautorefname}{বিভাগ}
\providecommand{\theoremautorefname}{উপপাদ্য}
\renewcommand{\theoremautorefname}{উপপাদ্য}
\end{code}
Existing custom reference names are preserved. Theorem environments that share
a counter also share hyperref's usual reference-type rules.

\subsection{Numeric citations}
For a short bibliography, you can write the entries directly and cite them as usual:
\begin{examplepair}
ইউক্লিডের বই দেখুন~\cite{euclid}।
\begin{thebibliography}{9}
\bibitem{euclid}
Euclid. \textit{Elements}.
\end{thebibliography}
\end{examplepair}
If you manage references with natbib or biblatex, the same numbering applies:
use natbib's numeric mode, including sorting and compression, or biblatex's
\pkg{numeric} or \pkg{numeric-comp} style. Only numeric labels become Bangla.
Labels such as \pkg{Knuth84}, years you type, notes and author--year citations
keep their script and style. Citation keys, sorting and links work as usual.
Set \pkg{changecounternumbering=false} to keep numeric labels unchanged too.

\subsection{A biblatex document}
For a longer article or book, keep your sources in a \pkg{.bib} file and let
biblatex build the bibliography and numbered citations.
Load LaTeXbangla before biblatex so Polyglossia is ready when biblatex starts.
Use this preamble and add \verb|\printbibliography| where the list belongs:
\begin{code}
\usepackage{latexbangla}
\setbanglamainfont{Kalpurush}
\usepackage[backend=biber,style=numeric-comp]{biblatex}
\addbibresource{references.bib}
\usepackage{hyperref}
\end{code}
Run your chosen LaTeX engine, Biber, then the same engine twice. The Bengali mapping translates headings
and conjunctions; remaining bibliography strings inherit English.
For natbib, choose a BibTeX style and load:
\begin{code}
\usepackage[numbers,sort&compress]{natbib}
\end{code}

\subsection{Using other packages}
Build on the packages you already use: \pkg{mathtools} and \pkg{bm} for
mathematics, or \pkg{booktabs}, \pkg{tabularx}, \pkg{longtable} and \pkg{multirow}
for tables. You can style captions with \pkg{caption} and \pkg{subcaption},
headers with \pkg{fancyhdr}, headings with \pkg{titlesec}, and boxed material
with \pkg{tcolorbox}. These combinations are tested; each package keeps control
of its own layout and styles.

For grouped or ranged references, load \pkg{cleveref} after \pkg{amsmath},
\pkg{amsthm} and \pkg{hyperref}. Counter and page numbers remain Bangla;
cleveref controls reference names and connecting words. For example:
\begin{code}
\usepackage[english]{cleveref}
\crefname{equation}{সমীকরণ}{সমীকরণসমূহ}
\AtBeginDocument{%
  \renewcommand{\crefrangeconjunction}{ থেকে }%
  \renewcommand{\crefpairconjunction}{ ও }%
}
\end{code}
Then use \verb|\cref{eq:first,eq:second}| or
\verb|\crefrange{eq:first}{eq:last}|. The \pkg{english} option supplies
cleveref's initial names and connecting words; the commands above customize them.
Load \pkg{microtype} normally to use its available typography features
with mixed Bangla and English fonts.

\section{Updating an older document}
You can keep using \pkg{banglamainfont} and \pkg{banglattfont} in existing
documents. When starting a new one, use the font commands shown in this guide;
you only need to choose one Bangla font. The table shows the equivalent settings.
\begin{center}
\begin{tabularx}{\linewidth}{@{}lX@{}}
\toprule
\textbf{Older setting} & \textbf{Current interface} \\
\midrule
\pkg{banglamainfont=Kalpurush} & \texttt{\detokenize{\setbanglamainfont{Kalpurush}}} \\
\pkg{banglattfont=FreeSerif} & \texttt{\detokenize{\setbanglamonofont{FreeSerif}}} \\
\pkg{changecounternumbering=1} & \pkg{changecounternumbering=true} (default) \\
\pkg{changecounternumbering=0} & \pkg{changecounternumbering=false} \\
\bottomrule
\end{tabularx}
\end{center}
A later font setter overrides its corresponding package option. The numbering
option remains current; both its Boolean and numeric spellings are accepted.
The old \pkg{SiyamRupali} font spelling is accepted as \pkg{Siyam Rupali}.

If your document uses \pkg{feature=0/1/2}, remove that option: language support
is now available without choosing a feature tier. Load the AMS, reference and
bibliography packages you need, and define your theorem or solution environments.
Margins, line spacing, heading
punctuation, theorem styles and equation reset rules belong to your document
class or style. Nested lists retain their decimal, alphabetic or Roman structure.

The old \pkg{latin} wrapper is replaced by automatic Bangla/English font
switching. Use Polyglossia's language commands when you want an actual language
change. Use \texttt{\detokenize{\eqref}} for a parenthesized equation number and
\texttt{\detokenize{\autoref}} for a named reference. Current Polyglossia supplies the corrected
digit converter, so the old numbering workaround is no longer needed.

\section*{Quick reference}
\phantomsection
\addcontentsline{toc}{section}{Quick reference}
\begin{center}
\begin{tabularx}{\linewidth}{@{}lX@{}}
\toprule
\textbf{Task} & \textbf{Command or setting} \\
\midrule
Choose the Bangla main font & \texttt{\detokenize{\setbanglamainfont[features]{font}}} \\
Choose optional Bangla roles & \texttt{\detokenize{\setbanglasansfont}}, \texttt{\detokenize{\setbanglamonofont}} \\
Choose English families & \texttt{\detokenize{\setmainfont}}, \texttt{\detokenize{\setsansfont}}, \texttt{\detokenize{\setmonofont}} \\
Keep existing counter formats & \pkg{changecounternumbering=false} \\
Print a counter in Bangla & \texttt{\detokenize{\tobangla{counter}}} \\
Convert a string of digits & \texttt{\detokenize{\numtobangla{0012}}} \\
Put Bangla words in math & \texttt{\detokenize{\text{বাংলা}}} \\
Alphabetic counter format & \texttt{\detokenize{\alph{counter}}} or \texttt{\detokenize{\Alph{counter}}} \\
\bottomrule
\end{tabularx}
\end{center}

With conversion enabled, alphabetic labels run from ক to ৎ (36 letters).
Zero is empty; negative values or values above 36 give ``Counter too large''.
Appendices use labels such as পরিশিষ্ট ক.

\subsection{Examples and help}
To see the features together, try the example document and font guide in
\href{https://github.com/Neehan/LaTeXbangla}{the repository},
both supplied with sources and PDFs. Tested classes include article, report, book and memoir.
If something does not work as expected, send a small document, your font names
and the compiler name and log; \texttt{\detokenize{\listfiles}} records package versions. You can reach us through
\href{https://github.com/Neehan/LaTeXbangla/issues}{GitHub Issues} or
\href{https://groups.google.com/g/latexbangla}{\texttt{latexbangla@googlegroups.com}}.

\smallskip
{\small Copyright 2016, 2026 Adib Hasan. Distributed under the
\href{https://www.latex-project.org/lppl.txt}{LaTeX Project Public License 1.3c or later}.
This manual uses Kalpurush for Bangla and TeX Gyre Pagella for English prose;
the opening example uses TeX Gyre Adventor. The family comparison also uses
Noto Sans Bengali and Siyam Rupali.}

\clearpage
% The translated guide uses the normal Bangla counter formatters. Restart its
% section and demonstration counters, but keep page numbers continuous. Give
% restarted counters distinct PDF anchors so links cannot jump to English pages.
\renewcommand{\thesection}{\arabic{section}}
\renewcommand{\thesubsection}{\thesection.\arabic{subsection}}
\renewcommand{\thepage}{\arabic{page}}
\renewcommand{\theHsection}{bangla.\number\value{section}}
\renewcommand{\theHequation}{bangla.\number\value{equation}}
\renewcommand{\theHtheorem}{bangla.\number\value{theorem}}
\setcounter{section}{0}
\setcounter{equation}{0}
\setcounter{theorem}{0}
\renewcommand{\manualcodetitle}{কোড}
\renewcommand{\manualresulttitle}{ফলাফল}
\fancyhead[R]{\small\sffamily\color{BanglaInk}সংস্করণ ১.১\quad ব্যবহার নির্দেশিকা}
\phantomsection\label{manual:bangla}
\addcontentsline{toc}{part}{বাংলা ব্যবহার নির্দেশিকা}
\thispagestyle{empty}
\begin{center}
\includegraphics[width=0.66\linewidth]{logo.png}\\[2pt]
{\Large\sffamily\color{BanglaGreen}চেনা \LaTeX-এর কমান্ডেই বাংলা লেখা}\\[5pt]
{আদীব হাসান}\\[2pt]
{\small\href{https://groups.google.com/g/latexbangla}{\texttt{latexbangla@googlegroups.com}}}\\[3pt]
{২৩ অক্টোবর ২০২৬\quad\textcolor{BanglaRed}{\textbullet}\quad সংস্করণ ১.১}
\end{center}

LaTeXbangla হলো পরিচিত \LaTeX-এর কমান্ড দিয়ে বাংলা প্রবন্ধ, প্রতিবেদন ও
বই লেখার একটি প্যাকেজ। শুরু করতে প্যাকেজটি যোগ করে একটি বাংলা ফন্ট বেছে নিন।
এরপর একই সঙ্গে বাংলা ও ইংরেজি লিখতে, \verb|\textbf| দিয়ে লেখা গাঢ় করতে,
আর তালিকা, সমীকরণ ও রেফারেন্স তৈরি করতে পারবেন আগের মতোই।
বাংলা ও ইংরেজির জন্য উপযুক্ত ফন্ট প্যাকেজ নিজেই বেছে নেয়। পৃষ্ঠার বিন্যাস,
ইংরেজি ফন্ট ও গণিতের ফন্ট নিজের পছন্দমতো রাখতে পারেন।

LaTeXbangla ডকুমেন্টের অন্য অংশগুলোও
বাংলায় সাজানোর ব্যবস্থা করে। যেমন তালিকায় ১ ও (ক)-এর মতো নম্বর আসে,
স্বয়ংক্রিয় রেফারেন্সে সমীকরণ-এর মতো বাংলা নাম আসে, আর সংখ্যাভিত্তিক
তথ্যসূত্রের নম্বরও বাংলায় দেখা যায়। প্যাকেজটি fontspec ও Polyglossia-এর ওপর
ভিত্তি করে তৈরি। এর বাড়তি সুবিধার মধ্যে আছে স্বয়ংক্রিয় ফন্ট পরিবর্তন,
গাঢ় ও তির্যক রূপের বিকল্প ব্যবস্থা, বাংলা বর্ণে ক্রমসংখ্যা, গণিতে বাংলা অঙ্ক
ও রেফারেন্সের বাংলা রূপ। পরের অংশগুলোতে উদাহরণসহ এগুলোর ব্যবহার দেখানো হয়েছে।

\section{দ্রুত আরম্ভ}
TeX Live 2026 বা পরের সামঞ্জস্যপূর্ণ সংস্করণে \textbf{XeLaTeX অথবা LuaLaTeX}
দিয়ে কম্পাইল করুন।
২ নম্বর পরিচ্ছেদের তালিকা থেকে
একটি বাংলা ফন্ট ইনস্টল করে এভাবে শুরু করুন:
\begin{tcolorbox}[manualbox,title=ন্যূনতম ডকুমেন্ট]
\begin{code}
\documentclass{article}
\usepackage{fontspec}
\setmainfont{TeX Gyre Adventor} % ঐচ্ছিক; ইংরেজি ফন্ট অক্ষুণ্ণ থাকে।
\usepackage{latexbangla}
\setbanglamainfont{Kalpurush}
\begin{document}
  % এখানে লিখুন।
\end{document}
\end{code}
\end{tcolorbox}
নিচের কোডটি \verb|\begin{document}| ও \verb|\end{document}|-এর মাঝে বসান।
খেয়াল করুন, গাঢ় ও তির্যক লেখাতেও ইংরেজি Adventor-এর নিজস্ব sans serif
চেহারা ধরে রাখে। প্রতিটি ফলাফলের বাক্সে পাশের কোডটিই চালানো হয়েছে।
\begingroup
\let\rmfamilylatin\manualexampleenglishfont
% Reselect after replacing the companion so both engines capture this local setup.
\bengalifont
\begin{examplepair}
বাংলা and English একই অনুচ্ছেদে।

English keeps its own typeface.

\textbf{গুরুত্বপূর্ণ: read this.}
\textit{বাংলা ও English একসঙ্গে।}
\end{examplepair}
\endgroup

\subsection{শুধু fontspec দিয়ে একই অনুচ্ছেদ}
ইংরেজি ডকুমেন্টে অল্প কিছু বাংলা লিখতে fontspec-ই অনেক সময় যথেষ্ট।
উপরের অনুচ্ছেদটি একইভাবে দেখাতে নিচের সম্পূর্ণ ডকুমেন্টটি ব্যবহার করতে পারেন।
বাংলার জন্য আলাদা ফন্ট বেছে Kalpurush-এর গাঢ় ও তির্যক রূপের সেটিং দিন,
তারপর প্রতিটি বাংলা অংশ সেই ফন্টের গ্রুপে রাখুন। গ্রুপের বাইরে ইংরেজি Adventor-এই থাকে।
LuaLaTeX-এ বাংলা অক্ষরের যুক্তরূপ ও অবস্থান ঠিক রাখতে HarfBuzz-ও চালু করুন:
\begin{tcolorbox}[manualbox,title=LaTeXbangla ছাড়া]
\begin{code}
\documentclass{article}
\usepackage{fontspec,iftex}
\setmainfont{TeX Gyre Adventor}
\ifluatex\defaultfontfeatures{Renderer=HarfBuzz}\fi
\newfontfamily\banglatext[
  Script=Bengali,
  AutoFakeBold=4,
  AutoFakeSlant=0.2
]{Kalpurush}
\begin{document}
{\banglatext বাংলা} and English
{\banglatext একই অনুচ্ছেদে।}

English keeps its own typeface.

\textbf{{\banglatext গুরুত্বপূর্ণ:} read this.}
\textit{{\banglatext বাংলা ও} English
  {\banglatext একসঙ্গে।}}
\end{document}
\end{code}
\end{tcolorbox}
LaTeXbangla ফন্ট পরিবর্তন ও রূপের এই সেটিংগুলো নিজেই সামলায়, তাই মিশ্র
অনুচ্ছেদটি সরাসরি লিখতে পারেন। উপরের উদাহরণে শুধু লেখার চেহারা মেলানো হয়েছে;
পরের অংশগুলোতে দেখানো বাংলা ক্রমসংখ্যা ও ডকুমেন্টের অনূদিত নামগুলো
LaTeXbangla-এর ভাষা সহায়তার আরও কিছু সুবিধা।

\textbf{স্থানীয়ভাবে ইনস্টল।} \href{https://github.com/Neehan/LaTeXbangla}{রিপোজিটরির}
\pkg{latexbangla/} ফোল্ডার খুলুন।\\
\pkg{latexbangla.sty}, \pkg{latexbangla-bengali.lbx} ও পুরো \pkg{code/}
ফোল্ডার মূল ডকুমেন্টের পাশে রাখুন; TeX tree-তেও একই কাঠামো ইনস্টল করা যায়।
পরীক্ষিত সংস্করণগুলো হলো LaTeX 2026-06-01, Polyglossia 2.13 ও fontspec 2.9h।

\section{ফন্ট ও লেখার ধরন বেছে নেওয়া}
\subsection{পরীক্ষিত বাংলা ফন্ট}
আপনার লেখার সঙ্গে মানায় এমন একটি ফন্ট দিয়ে শুরু করুন। নিচের যেকোনো
পরীক্ষিত ফন্ট দিয়েই পুরো বাংলা ডকুমেন্ট লেখা যায়; নামটি বসান
\verb|\setbanglamainfont{...}|-এ। Anek-এর জন্য নিচের নির্দেশনা অনুযায়ী ফাইলের নাম দিন। বেছে নেওয়ার আগে
\href{https://github.com/Neehan/LaTeXbangla}{রিপোজিটরিতে দেওয়া}
ফন্ট তুলনার PDF-এ নমুনাগুলো দেখে নিতে পারেন।
\begin{center}
\begin{tabularx}{\linewidth}{@{}lX@{}}
\toprule
\textbf{ফন্টের নাম} & \textbf{গাঢ় বাংলা লেখার স্বাভাবিক সেটিং} \\
\midrule
\texttt{Kalpurush} & কৃত্রিম গাঢ়, মাত্রা ৪ \\
\texttt{SolaimanLipi} & কৃত্রিম গাঢ়, মাত্রা ৪ \\
\texttt{Siyam Rupali} & কৃত্রিম গাঢ়, মাত্রা ৪.৫ \\
\texttt{Noto Serif Bengali} & ফন্টের নিজস্ব গাঢ় রূপ \\
\texttt{Noto Sans Bengali} & ফন্টের নিজস্ব গাঢ় রূপ \\
\texttt{Anek Bangla} & ভ্যারিয়েবল ফাইলে কৃত্রিম গাঢ়, মাত্রা ৪ \\
\texttt{Tiro Bangla} & কৃত্রিম গাঢ়, মাত্রা ৩.৫; নিজস্ব তির্যক রূপ \\
\texttt{Galada} & কৃত্রিম গাঢ়, মাত্রা ৩.৫ \\
\texttt{FreeSerif} & কৃত্রিম গাঢ়, মাত্রা ৩.৫; নিজস্ব তির্যক রূপ \\
\bottomrule
\end{tabularx}
\end{center}
ফন্টের নিজস্ব গাঢ় বা তির্যক রূপে বাংলা অক্ষর থাকলে LaTeXbangla সেটিই নেয়।
না থাকলে অক্ষরের রেখা মোটা করে বা হেলিয়ে সেই রূপ তৈরি করে। তালিকার বাইরের
ফন্টে কৃত্রিম রূপ দরকার হলে গাঢ় করার মাত্রা ২.৫ ও তির্যক করার মাত্রা ০.২। Small caps ও অন্য রূপের
প্রাপ্যতা ফন্টের ওপর নির্ভর করে। ফন্ট আলাদা করে ইনস্টল করুন।

\begin{tcolorbox}[manualwarning]
Anek Bangla নাম দিয়ে লোড না হলে TTF ফাইলটি \texttt{AnekBangla-VF.ttf} নামে
ডকুমেন্টের পাশে রেখে \verb|\setbanglamainfont{AnekBangla-VF.ttf}| লিখুন।
\end{tcolorbox}

\subsection{বাংলা ও ইংরেজির আলাদা ফন্ট}
বাংলা ও ইংরেজি একসঙ্গে রেখেই মূল লেখা, শিরোনাম ও টাইপরাইটার ধরনের লেখায়
আলাদা চেহারা দিতে পারেন। এখানে মূল লেখার জন্য Kalpurush ও Pagella,
sans serif-এর জন্য Noto Sans Bengali ও Heros, আর mono-র জন্য Siyam Rupali
ও Cursor বেছে নেওয়া হয়েছে। সেটিংগুলো preamble-এ, অর্থাৎ
\verb|\begin{document}|-এর আগে লিখুন:
\begin{tcolorbox}[manualbox,title=আলাদা ফন্ট বাছার উদাহরণ]
\begin{code}
\setbanglamainfont{Kalpurush}
\setbanglasansfont{Noto Sans Bengali}
\setbanglamonofont{Siyam Rupali}
\setmainfont{TeX Gyre Pagella}
\setsansfont{TeX Gyre Heros}
\setmonofont{TeX Gyre Cursor}
\end{code}
\end{tcolorbox}
এবার আগের মতোই \verb|\textrm|, \verb|\textsf| ও \verb|\texttt| ব্যবহার করুন।
প্রতিটি কমান্ড সংশ্লিষ্ট ফন্টের জোড়া বেছে নেয়; গাঢ় ও তির্যক রূপও কাজ করে:
\begingroup
\let\bengalifontsf\manualsansfont
\let\bengalifonttt\manualmonofont
\begin{examplepair}
বাংলা ও English।\\
\textbf{গাঢ় ও bold।}\\
\textit{তির্যক ও italic।}

\textsf{বাংলা ও English।}\\
\textsf{\textbf{গাঢ় ও bold।}}\\
\textsf{\textit{তির্যক ও italic।}}

\texttt{বাংলা ও English।}\\
\texttt{\textbf{গাঢ় ও bold।}}\\
\texttt{\textit{তির্যক ও italic।}}
\end{examplepair}
\endgroup
বাংলা বা ইংরেজি অংশ আলাদা করে চিহ্নিত না করেই একই লেখা তিন রকম দেখায়।
Siyam Rupali-র অক্ষরের স্বাভাবিক প্রস্থ বজায় থাকে; mono-র জন্য বাছলেই
ফন্টটির সব অক্ষর সমান চওড়া হয়ে যায় না।

এই আলাদা সেটিংগুলো ঐচ্ছিক: একটি মূল বাংলা ফন্ট দিয়েই তিন ধরনের লেখা চলে,
আর ইংরেজির ফন্ট না বদলালে আগের ফন্টই থাকে। বাংলা ও ইংরেজির কমান্ড যেকোনো
ক্রমে লেখা যায়; প্রতিটি কমান্ড শুধু তার জন্য নির্দিষ্ট ফন্টটি বদলায়।
স্বয়ংক্রিয় ফন্ট পরিবর্তন লেখার ভাষা বা লাইনের শেষে হাইফেন দিয়ে শব্দ ভাঙার
নিয়ম বদলায় না। ইংরেজি অংশে ইংরেজির নিয়ম ব্যবহার করতে preamble-এ
\verb|\setotherlanguage{english}| দিন, তারপর \verb|\textenglish{...}| ব্যবহার করুন।

\subsection{Fontspec-এর অপশন, ফন্ট ফাইল ও ফাঁক}
অন্য ফন্টের সঙ্গে আকার মেলাতে, শব্দের ফাঁক বদলাতে বা নিজের প্রজেক্টের ফন্ট
ফাইল ব্যবহার করতে বাংলা ফন্ট কমান্ডে fontspec-এর অপশন দিন। তিনটি কমান্ডেই
ফন্টের নামের আগে বা পরে অপশন লেখা যায়:
\begin{code}
\setbanglamainfont[features]{font}
\setbanglamainfont{font}[features]
\end{code}
এগুলো fontspec-এরই অপশন; \pkg{Script=Bengali} আগে থেকে দেওয়া থাকে।
LuaLaTeX-এ স্বাভাবিকভাবে \pkg{Renderer=HarfBuzz}-ও দেওয়া থাকে।
আপনার দেওয়া মান প্যাকেজের স্বাভাবিক সেটিংকে প্রতিস্থাপন করে। যেমন আকার
ও শব্দের ফাঁক বদলাতে বা ফাইল থেকে ফন্ট নিতে লিখতে পারেন:
\begin{code}
\setbanglamainfont[Scale=1.05,WordSpace=1.1]{Kalpurush}
\setbanglamainfont[Path=fonts/]{fontkalpurush.ttf}
\end{code}
\pkg{WordSpace=1.1} ফন্টের নিজস্ব শব্দের ফাঁক ১০\% বাড়ায়;
প্যাকেজের পূর্বনির্ধারিত গুণকের বদলে এই মান ব্যবহৃত হয়। আকার ও ফাঁকের
স্বাভাবিক সেটিংগুলো হলো:
\begin{center}
\begin{tabular}{lcc}
\toprule
ফন্ট & Scale & WordSpace \\
\midrule
Kalpurush & ১ & ১ \\
SolaimanLipi & ১ & ১.২ \\
Siyam Rupali & ০.৮ & ১.১২ \\
Noto Serif Bengali & ১ & ১.৩ \\
Noto Sans Bengali & ১ & ১.২ \\
Anek Bangla & ০.৯ & ২.১ \\
Tiro Bangla & ১ & ১.২৫ \\
Galada & ০.৯ & ১.৬৭ \\
FreeSerif & ১ & ১.৩৩ \\
\bottomrule
\end{tabular}
\end{center}
পরিচিত নাম ও প্রচলিত ফাইলের নাম অনুযায়ী সেটিং চেনা হয়;
ফাইলের নাম বদলালে সাধারণ সেটিং প্রযোজ্য।
আলাদা ফাইলে রাখা রূপ বাছতে \pkg{UprightFont}, \pkg{BoldFont},
\pkg{ItalicFont} ও \pkg{BoldItalicFont} ব্যবহার করুন।

\section{গাঢ় ও তির্যক লেখা নিয়ন্ত্রণ}
কোনো সংজ্ঞা চোখে পড়ার মতো করতে \verb|\textbf|, আর জোর দিয়ে কিছু বোঝাতে
\verb|\textit| ব্যবহার করুন; সাধারণ সেটিংয়েই দুটো কাজ করে। আপনার ডকুমেন্টে
কৃত্রিম গাঢ় লেখা আরেকটু স্পষ্ট চাইলে \pkg{AutoFakeBold}-এর মান বাড়াতে পারেন।
গাঢ়-তির্যক লেখাতেও একই মাত্রা রাখতে সঙ্গে \pkg{AutoFakeSlant} দিন।
নিচের উদাহরণে গাঢ় করার মাত্রা ৫ দিয়ে চারটি রূপ দেখানো হয়েছে:
\begin{tcolorbox}[manualbox,title=গাঢ় ও তির্যক রূপের সেটিং (preamble)]
\begin{code}
\setbanglamainfont[
  AutoFakeBold=5,
  AutoFakeSlant=0.2
]{Kalpurush}
\end{code}
\end{tcolorbox}
\begingroup
\let\bengalifont\manualstrongfont
\bengalifont
\begin{examplepair}
সাধারণ লেখা।

\textbf{গাঢ় লেখা।}

\textit{তির্যক লেখা।}

\textbf{\textit{গাঢ় ও তির্যক।}}
\end{examplepair}
\endgroup
নিজে \pkg{AutoFakeBold} দিলে ফন্টের নিজস্ব গাঢ় রূপ থাকলেও কৃত্রিম গাঢ়
ব্যবহৃত হবে। তাই সংখ্যাসহ অপশনটি দিন: শুধু \pkg{AutoFakeBold} লিখলে
fontspec ১.৫ ধরে, LaTeXbangla-এর ফন্ট অনুযায়ী নির্ধারিত মান ধরে না।

গাঢ় ও গাঢ়-তির্যক রূপ আলাদাভাবে বদলাতেও fontspec-এর অপশন ব্যবহার করতে পারেন:
\begin{code}
\setbanglamainfont[
  BoldFeatures={FakeBold=5},
  BoldItalicFeatures={FakeBold=5,FakeSlant=0.2}
]{Kalpurush}
\end{code}
কৃত্রিম রূপ দরকার এমন ফন্টে এই অপশনগুলো বিশেষ কাজে আসে। পরীক্ষিত সেটিং
ও নিজস্ব রূপের স্বয়ংক্রিয় নির্বাচন চাইলে সাধারণ ফন্ট কমান্ডই যথেষ্ট।
Polyglossia-এর \verb|\bengalifont|, \verb|\bengalifontsf| ও
\verb|\bengalifonttt| সরাসরি সংজ্ঞায়িত করেও ফন্ট বাছা যায়;
সেক্ষেত্রে রূপগুলোর সেটিং নিজে দিন। আগে থেকে সংজ্ঞায়িত কমান্ড বদলাতে
\verb|\renewfontfamily| ব্যবহার করুন; LaTeXbangla-এর দেওয়া sans ও mono
ফন্টের কমান্ডও এর মধ্যে পড়ে।

\subsection{যতিচিহ্ন, তালিকার বুলেট ও পাদটীকা}
মিশ্র লেখাতেও যতিচিহ্ন ও পাদটীকা আগের মতোই ব্যবহার করতে পারেন। সরাসরি লেখা ড্যাশ,
বাঁকানো উদ্ধৃতিচিহ্ন ও তিন-বিন্দু চিহ্ন (…) ইংরেজি ফন্ট থেকে আসে। বাংলা ফন্টে চিহ্ন না
থাকলে সাধারণ বুলেট ও প্রতীকভিত্তিক পাদটীকার চিহ্নও সেখান থেকে নেওয়া হয়।
সাধারণ পাদটীকায় বাংলা নম্বর আসে; Roman বা প্রতীকভিত্তিক বিন্যাস বাছলে তা-ই থাকে।
\begin{examplepair}
“বাংলা”—English… আবার বাংলা।
\begin{itemize}
  \item একটি কথা।
  \item আরেকটি কথা।
\end{itemize}
একটি তথ্য\footnote{ছোট ব্যাখ্যা।}
\end{examplepair}

\section{ক্রমসংখ্যা, তালিকা ও নিজের কাউন্টার}\label{bn:sec:numbering}
উপপ্রশ্নসহ একটি অনুশীলনী লিখতে সাধারণ তালিকার ভেতরে আরেকটি তালিকা ব্যবহার করুন।
বাংলা নম্বর LaTeXbangla নিজেই দেয়, তাই পুরো ডকুমেন্টে পরিচিত তালিকার কমান্ডই চলে।

স্বাভাবিক সেটিং হলো \pkg{changecounternumbering=true}: অঙ্কে লেখা কাউন্টারের
নম্বর বাংলা হয়, আর \verb|\alph| ও \verb|\Alph|-এ বাংলা বর্ণ আসে। কাউন্টারের মান,
পুনরায় গণনা শুরুর নিয়ম ও যৌগিক নম্বরের কাঠামো একই থাকে। Roman ও
প্রতীকভিত্তিক বিন্যাসও বদলায় না। নির্দেশিকার এই বাংলা অংশে পরিচ্ছেদ ও
পাতার নম্বরসহ সব সাধারণ কাউন্টারেই বাংলা ব্যবহার করা হয়েছে। সাধারণ
\pkg{enumerate}-এর চারটি স্তরে যথাক্রমে ১., (ক), i. ও ক. আসে; যতিচিহ্ন
ও Roman সংখ্যা সাধারণ \LaTeX-এর মতোই থাকে।
\begin{examplepair}
\begin{enumerate}
  \item একটি আয়তক্ষেত্রের দৈর্ঘ্য ৪,
    প্রস্থ ৩ একক।
  \begin{enumerate}
    \item ক্ষেত্রফল কত?
    \item কর্ণের দৈর্ঘ্য কত?
  \end{enumerate}
\end{enumerate}
\end{examplepair}

\subsection{নিজের কাউন্টারও সরাসরি কাজ করে}
অনুশীলনীর নম্বর দিতে আগের মতোই একটি কাউন্টার তৈরি করুন।
\verb|\refstepcounter| মান এক বাড়ায় এবং \verb|\label| ও \verb|\ref|-এর
ব্যবস্থা করে; \verb|\stepcounter| শুধু মান বাড়ায়।
\begin{examplepair}
\newcounter{practice}
\setcounter{practice}{11}
\refstepcounter{practice}
\label{bn:ex:custom}
অনুশীলনী \thepractice।\par
\renewcommand{\thepractice}
  {\alph{practice}}
\setcounter{practice}{2}
অনুশীলনী \thepractice।
\end{examplepair}
এখানে \texttt{[section]} প্রতি পরিচ্ছেদে কাউন্টার শূন্য করে;
\verb|\theproblem|-এর সংজ্ঞা মুদ্রিত নম্বরে পরিচ্ছেদের নম্বর যোগ করে:
\begin{code}
\newcounter{problem}[section]
\renewcommand{\theproblem}{\thesection.\arabic{problem}}
\end{code}
চাইলে শুধু একটি কাউন্টারের জন্যও বাংলা বা ইংরেজি অঙ্ক বেছে নিতে পারেন:
\begin{code}
\renewcommand{\thepractice}{\tobangla{practice}}       % বাংলা
\renewcommand{\thepractice}{\number\value{practice}} % ইংরেজি
\end{code}
\verb|\arabic|, \verb|\alph| বা \verb|\Alph| ব্যবহার না করে লেখা
বিন্যাস নিজের ফলাফল অপরিবর্তিত রাখে।

\subsection{নির্দিষ্ট কাউন্টারের নম্বর ইংরেজিতে রাখা}
ডকুমেন্টের বাকি নম্বর বাংলায় রেখেও একটি কাউন্টারের নম্বর ইংরেজিতে
লিখতে পারেন। সেই কাউন্টারের রেফারেন্সেও একই বিন্যাস থাকে:
\begin{examplepair}
\renewcommand{\thepractice}
  {\number\value{practice}}
\refstepcounter{practice}
\label{bn:ex:western}
অনুশীলনী \thepractice।\par
অনুশীলনী~\ref{bn:ex:western} দেখুন।
\end{examplepair}
ইংরেজি বর্ণের জন্য প্রিঅ্যাম্বলে \verb|\usepackage{alphalph}| লিখুন।
এর \verb|\alphalph| বিন্যাসে a, b, \ldots, z, aa, \ldots\ আসে;
\verb|\AlphAlph| বড় হাতের বর্ণ দেয়। নিচে সংজ্ঞাটি তালিকার পরিবেশের
ভেতরে থাকায় শুধু এই তালিকার নম্বর বদলায়:
\begin{examplepair}
\begin{enumerate}
  \renewcommand{\theenumi}
    {\alphalph{\value{enumi}}}
  \item প্রথম।\label{bn:item:latin}
  \item দ্বিতীয়।
\end{enumerate}
দফা~\ref{bn:item:latin} দেখুন।
\end{examplepair}
এই তালিকায় 1, 2, \ldots\ চাইলে বদলে \verb|\number\value{enumi}| লিখুন।
নিজের কাউন্টারের জন্য \verb|\thepractice|-এর সংজ্ঞায়
\verb|\alphalph{\value{practice}}| দিন। পরের তালিকাগুলোর নম্বর বাংলাই থাকে।
এই পরিবর্তন শুধু নম্বরের জন্য; \verb|\autoref|-এর দেওয়া অনূদিত নাম বদলায় না।

\subsection{একাধিক বর্ণের নম্বর}
সাধারণ বর্ণভিত্তিক নম্বর ৩৬তম বর্ণ ৎ-এ শেষ হয়। আরও নম্বর দরকার হলে
\pkg{alphalph} দিয়ে এর পরে কক, কখ ইত্যাদি লেখা যায়। কাউন্টার রূপান্তর
চালু রেখে, LaTeXbangla লোড করার পরে প্রিঅ্যাম্বলে লিখুন:
\begin{code}
\usepackage{alphalph}
\makeatletter
\newalphalph{\extendedBangla}{\@alph}{36}
\makeatother
\end{code}
শেষ আর্গুমেন্টটি বর্ণের সংখ্যা। যে কাউন্টারে একাধিক বর্ণের নম্বর চান,
সেটিতে নতুন বিন্যাসটি ব্যবহার করুন; অন্য কাউন্টারের আচরণ বদলাবে না:
\begin{examplepair}
\newcounter{extendedpractice}
\renewcommand{\theextendedpractice}
  {\extendedBangla{
    \value{extendedpractice}}}
\setcounter{extendedpractice}{36}
\theextendedpractice,
\stepcounter{extendedpractice}
\theextendedpractice,
\stepcounter{extendedpractice}
\theextendedpractice
\end{examplepair}
\verb|\label| ও \verb|\ref| ব্যবহার করতে চাইলে \verb|\stepcounter|-এর
বদলে \verb|\refstepcounter| লিখুন। প্যাকেজের নিজস্ব \verb|\alphalph| ও
\verb|\AlphAlph| কমান্ডে ইংরেজি বর্ণই আসে; বাংলা নম্বরের জন্য এখানে
তৈরি করা \verb|\extendedBangla| কমান্ডটি ব্যবহার করুন।

\subsection{স্বয়ংক্রিয় বা সরাসরি রূপান্তর}
\verb|\usepackage[changecounternumbering=false]{latexbangla}| দিলে
কাউন্টার ও সংখ্যাভিত্তিক তথ্যসূত্রের আগের বিন্যাস থাকে; ফন্ট ও অনুবাদ কাজ করে।
এটি প্যাকেজের অপশন; ভাষা বদলালে এর মান বদলায় না। শুরুর শূন্যসহ কোনো কোড
লিখতে, বা স্বয়ংক্রিয় রূপান্তর বন্ধ রেখেও কাউন্টারের মান বাংলায় দেখাতে লিখুন:
\begin{examplepair}
\numtobangla{0012}; \tobangla{practice}
\end{examplepair}
কমান্ড দুটি expandable। \verb|\numtobangla| শুধু ইংরেজি অঙ্ক নেয় এবং শুরুর শূন্যগুলো রাখে;
ধনাত্মক-ঋণাত্মক চিহ্ন বা দশমিক বিন্দু নেয় না। \verb|\tobangla|
কাউন্টারের পূর্ণসংখ্যা মান বাংলায় লেখে; শূন্য ও ঋণাত্মক মানও চলে।

\section{গণিত, সারণি ও উপপাদ্য}
বাংলা পাঠে সূত্রের ব্যাখ্যা বাংলায় লিখতে, $x$ ও $v$-এর মতো চেনা প্রতীক
রাখতে এবং দুই লিপিতেই সংখ্যা লিখতে পারেন।
সংখ্যা যে লিপিতে দেখতে চান, সেই লিপিতেই লিখুন: \texttt{12} থাকে 12,
\texttt{১২} থাকে ১২, লেখায় ও গণিতে দুজায়গাতেই। কাউন্টারের স্বয়ংক্রিয়
রূপান্তর আপনার সূত্রের অঙ্ক বদলায় না। গণিতে বাংলা অঙ্ক মূল বাংলা ফন্ট থেকে আসে;
সূত্রে তার অবস্থান অনুযায়ী আকার মেলে। অন্য চিহ্নগুলো আসে গণিতের ফন্ট থেকে।

\pkg{amsmath} যোগ করে \verb|\text| দিয়ে রাশি ও এককের নাম বাংলায় লিখুন।
১২ মিটার পথ যেতে ৩ সেকেন্ড লাগলে লেখা যায়:
\begin{examplepair}
\begin{equation}\label{bn:eq:speed}
  \text{বেগ}=
  \frac{১২}{৩}=৪\text{ মি./সে.}
\end{equation}
\[
  x_৩^{২}+\sqrt{৪}=12+১২
\]
\end{examplepair}
সমীকরণের নম্বরও বাংলায় আসে; পরে \verb|\eqref| দিয়ে সেটি উল্লেখ করতে পারেন।
একইভাবে ভগ্নাংশ, মূল, ম্যাট্রিক্স ও একাধিক স্তরের সূচকেও বাংলা অঙ্ক লিখতে পারেন।
একটি অঙ্ক হলে \verb|x^২|-এর মতো বন্ধনী ছাড়াই চলে; একাধিক অঙ্কে বন্ধনী দিন।
গণিত গাঢ় করতে \verb|\boldmath| ব্যবহার করুন; শুধু বাইরের \verb|\textbf|
গণিতের লেখা গাঢ় করে না। \verb|\mathbb|-এর মতো বিশেষ গাণিতিক বর্ণমালার কমান্ড
বাংলা অঙ্কের আলাদা নকশা বেছে দেয় না।
অন্য গাণিতিক চিহ্নের গাঢ় রূপ আপনার গণিতের ফন্টের সেটিংয়ের ওপর নির্ভর করে।

\subsection{গণিতের ফন্ট আলাদা করে বাছুন}
সূত্রের চেহারা বদলাতে \pkg{unicode-math} দিয়ে পছন্দের গণিতের ফন্ট বেছে নিন:
\begin{code}
\usepackage{unicode-math}
\setmathfont{TeX Gyre Pagella Math}
\end{code}
প্রচলিত পদ্ধতিতে গণিতের ফন্ট বাছতে \pkg{latexbangla}-এর \emph{আগে}
\pkg{amsmath,newtxmath} যোগ করুন। গণিতের ফন্ট নির্বাচন আপনার হাতেই থাকে।

\subsection{সারণি ও নিজের তৈরি পরিবেশ}
সাধারণ সারণিতে বাংলা বর্ণনার পাশে ইংরেজি শব্দ ও মাপ লিখুন।
প্রতিটি ঘরে স্বয়ংক্রিয়ভাবে উপযুক্ত ফন্ট ব্যবহৃত হয়:
\begin{examplepair}
\begin{tabular}{llc}
বাহু & Side & দৈর্ঘ্য \\
\hline
প্রথম & First & ৩ \\
দ্বিতীয় & Second & ৪
\end{tabular}
\end{examplepair}
\pkg{siunitx}-এর সংখ্যার ঘরে ইংরেজি অঙ্ক দিন;
সরাসরি বাংলা অঙ্ক লিখতে চাইলে সাধারণ লেখার কলাম ব্যবহার করুন।
উপপাদ্য সংজ্ঞায়িত করতে \pkg{amsthm} ব্যবহার করুন। যেমন preamble-এ লিখুন:
\begin{code}
\newtheorem{theorem}{উপপাদ্য}
\end{code}
প্রমাণের বাংলা শিরোনাম Polyglossia দেয়;
উপস্থাপনার নকশা ডকুমেন্ট ক্লাস বা theorem style অনুযায়ী হয়।
\begin{examplepair}
\begin{theorem}
\label{bn:thm:pythagoras}
সমকোণী ত্রিভুজের লম্ব দুই বাহুর
দৈর্ঘ্য ৩ ও ৪ হলে অতিভুজের দৈর্ঘ্য ৫।
\end{theorem}
\begin{proof}
$৩^{২}+৪^{২}=৫^{২}$।
\end{proof}
\end{examplepair}

\section{লেবেল, রেফারেন্স ও তথ্যসূত্র}
লেবেল দিয়ে সমীকরণ, উপপাদ্য বা অনুশীলনীর উল্লেখ করলে পরে ডকুমেন্ট বদলালেও নম্বর
নিজে থেকে ঠিক থাকে। এখানে পরিচিত রেফারেন্স কমান্ডগুলো দিয়েই আগের
উদাহরণগুলোর উল্লেখ করা হয়েছে।
অধ্যায়, চিত্র, সারণি, সূচিপত্র, পরিশিষ্ট ও তথ্যসূত্রের শিরোনাম এবং তারিখের
বাংলা রূপ আসে Polyglossia থেকে। \pkg{hyperref} যোগ করলে LaTeXbangla
স্বয়ংক্রিয় রেফারেন্সের নামও বাংলায় দেয়। নতুন লেবেল বা তথ্যসূত্র যোগ করলে আবার কম্পাইল করুন,
যাতে রেফারেন্সের নম্বরগুলো ঠিকভাবে আসে।
\begin{examplepair}
\autoref{bn:eq:speed};
\eqref{bn:eq:speed}।

\autoref{bn:thm:pythagoras}।
অনুশীলনী \ref{bn:ex:custom}।
\end{examplepair}
\verb|\ref| নম্বর লেখে, \verb|\eqref| তার চারপাশে বন্ধনী দেয়,
আর \verb|\autoref| অনূদিত নামসহ রেফারেন্স দেয়। নামগুলোর মধ্যে আছে
পরিচ্ছেদ, উপপরিচ্ছেদ, সমীকরণ ও পাদটীকা। অন্য নাম চাইলে প্যাকেজগুলো
লোড করার পরে hyperref-এর সংশ্লিষ্ট কমান্ড বদলান:
\begin{code}
\renewcommand{\sectionautorefname}{বিভাগ}
\providecommand{\theoremautorefname}{উপপাদ্য}
\renewcommand{\theoremautorefname}{উপপাদ্য}
\end{code}
আগে থেকে দেওয়া নিজস্ব রেফারেন্সের নাম অক্ষুণ্ণ থাকে। একাধিক উপপাদ্যের
পরিবেশ একই কাউন্টার ব্যবহার করলে রেফারেন্সের ধরনেও hyperref-এর সাধারণ নিয়ম প্রযোজ্য।

\subsection{সংখ্যাভিত্তিক তথ্যসূত্র}
তথ্যসূত্রের তালিকা ছোট হলে সরাসরি লিখে আগের মতোই \verb|\cite| দিয়ে উল্লেখ করতে পারেন:
\begin{examplepair}
ইউক্লিডের বই দেখুন~\cite{bn:euclid}।
\begin{thebibliography}{9}
\bibitem{bn:euclid}
Euclid. \textit{Elements}.
\end{thebibliography}
\end{examplepair}
natbib-এর numeric mode-এ sorting ও compression-সহ সংখ্যার বাংলা রূপান্তর কাজ করে।
biblatex-এ \pkg{numeric} ও \pkg{numeric-comp} style ব্যবহার করতে পারেন।
শুধু অঙ্ক দিয়ে তৈরি লেবেল বাংলা হয়; \pkg{Knuth84}-এর মতো লেখা,
নিজে লেখা সাল, টীকা ও লেখক-সালভিত্তিক তথ্যসূত্রের লিপি ও বিন্যাস একই থাকে।
তথ্যসূত্রের key, সাজানোর ক্রম ও লিংক আগের মতোই কাজ করে।
\pkg{changecounternumbering=false} দিলে সংখ্যাভিত্তিক লেবেলও আগের মতো থাকে।

\subsection{Biblatex দিয়ে ডকুমেন্ট}
দীর্ঘ প্রবন্ধ বা বইয়ের জন্য তথ্যসূত্রগুলো \pkg{.bib} ফাইলে রাখুন। biblatex
সেই ফাইল থেকে তালিকা ও সংখ্যাভিত্তিক তথ্যসূত্র তৈরি করতে পারে।
biblatex-এর আগে LaTeXbangla যোগ করুন, যাতে Polyglossia প্রস্তুত থাকে।
নিচের অংশটি preamble-এ দিন; তথ্যসূত্রের তালিকা যেখানে চান সেখানে
\verb|\printbibliography| লিখুন:
\begin{code}
\usepackage{latexbangla}
\setbanglamainfont{Kalpurush}
\usepackage[backend=biber,style=numeric-comp]{biblatex}
\addbibresource{references.bib}
\usepackage{hyperref}
\end{code}
আপনার নির্বাচিত LaTeX ইঞ্জিন, তারপর Biber, তারপর একই ইঞ্জিন আরও দুবার চালান। বাংলা mapping শিরোনাম ও
সংযোজক শব্দ অনুবাদ করে; বাকি শব্দ ইংরেজি mapping থেকে আসে।
natbib-এর জন্য একটি BibTeX style বেছে নিয়ে যোগ করুন:
\begin{code}
\usepackage[numbers,sort&compress]{natbib}
\end{code}

\subsection{অন্য প্যাকেজের সঙ্গে ব্যবহার}
পরিচিত প্যাকেজগুলো নিয়েই কাজ করতে পারেন: গণিতের জন্য \pkg{mathtools} ও
\pkg{bm}, আর সারণির জন্য \pkg{booktabs}, \pkg{tabularx}, \pkg{longtable} ও
\pkg{multirow}। ক্যাপশন সাজাতে \pkg{caption} ও \pkg{subcaption}, পাতার ওপরের
লেখা সাজাতে \pkg{fancyhdr}, শিরোনামের জন্য \pkg{titlesec}, আর বাক্সে লেখা
রাখতে \pkg{tcolorbox} ব্যবহার করতে পারেন। এগুলোর সঙ্গে প্যাকেজটি পরীক্ষা
করা হয়েছে; প্রতিটির বিন্যাস ও নকশা সেই প্যাকেজই নিয়ন্ত্রণ করে।

একসঙ্গে কয়েকটি রেফারেন্স বা একটি পরিসর লিখতে \pkg{amsmath}, \pkg{amsthm}
ও \pkg{hyperref}-এর পরে \pkg{cleveref} লোড করুন। কাউন্টার ও পাতার নম্বর
বাংলাই থাকে; রেফারেন্সের নাম ও সংযোজক শব্দ cleveref নিয়ন্ত্রণ করে। যেমন:
\begin{code}
\usepackage[english]{cleveref}
\crefname{equation}{সমীকরণ}{সমীকরণসমূহ}
\AtBeginDocument{%
  \renewcommand{\crefrangeconjunction}{ থেকে }%
  \renewcommand{\crefpairconjunction}{ ও }%
}
\end{code}
এরপর \verb|\cref{eq:first,eq:second}| বা
\verb|\crefrange{eq:first}{eq:last}| ব্যবহার করুন। \pkg{english} অপশনটি
cleveref-এর প্রাথমিক নাম ও সংযোজক শব্দ দেয়; ওপরের কমান্ডগুলো সেগুলো বদলায়।
বাংলা ও ইংরেজি ফন্ট একসঙ্গে ব্যবহার করেও \pkg{microtype} স্বাভাবিকভাবে লোড
করে আপনার ইঞ্জিনে এর উপলব্ধ টাইপোগ্রাফির সুবিধাগুলো ব্যবহার করতে পারেন।

\section{পুরোনো ডকুমেন্ট হালনাগাদ}
পুরোনো ডকুমেন্টে \pkg{banglamainfont} ও \pkg{banglattfont} অপশন রেখেই কাজ
চালিয়ে যেতে পারেন। নতুন ডকুমেন্টে এই নির্দেশিকার ফন্ট কমান্ডগুলো ব্যবহার করুন;
একটি বাংলা ফন্ট বাছলেই চলে। নিচের ছকে পুরোনো সেটিংয়ের সমতুল্য নতুন কমান্ড আছে।
\begin{center}
\begin{tabularx}{\linewidth}{@{}lX@{}}
\toprule
\textbf{পুরোনো সেটিং} & \textbf{বর্তমান কমান্ড বা অপশন} \\
\midrule
\pkg{banglamainfont=Kalpurush} & \texttt{\detokenize{\setbanglamainfont{Kalpurush}}} \\
\pkg{banglattfont=FreeSerif} & \texttt{\detokenize{\setbanglamonofont{FreeSerif}}} \\
\pkg{changecounternumbering=1} & \pkg{changecounternumbering=true} (স্বাভাবিক সেটিং) \\
\pkg{changecounternumbering=0} & \pkg{changecounternumbering=false} \\
\bottomrule
\end{tabularx}
\end{center}
পরে ফন্ট বাছার কমান্ড দিলে সেটি সংশ্লিষ্ট প্যাকেজ অপশনকে প্রতিস্থাপন করে।
ক্রমসংখ্যার অপশনটি বর্তমান ইন্টারফেসেরই অংশ; Boolean ও সংখ্যায় লেখা দুই
রূপই গ্রহণ করা হয়। পুরোনো \pkg{SiyamRupali} নামটিও \pkg{Siyam Rupali} হিসেবে চলে।

ডকুমেন্টে \pkg{feature=0/1/2} থাকলে সেটি বাদ দিন; ভাষার সুবিধা পেতে এখন
আলাদা স্তর বাছতে হয় না। প্রয়োজনীয় AMS, রেফারেন্স ও তথ্যসূত্রের প্যাকেজ যোগ করে নিজের
উপপাদ্য বা সমাধানের পরিবেশ তৈরি করুন। মার্জিন, লাইনের ফাঁক, শিরোনামের
যতিচিহ্ন, উপপাদ্যের নকশা ও সমীকরণের গণনা পুনরায় শুরুর নিয়ম ডকুমেন্ট ক্লাস
বা style-এর বিষয়। একাধিক স্তরের তালিকায় অঙ্ক, বর্ণ বা Roman সংখ্যায় নম্বর দেওয়ার কাঠামো বজায় থাকে।

পুরোনো \pkg{latin} পরিবেশের কাজ এখন স্বয়ংক্রিয় বাংলা-ইংরেজি ফন্ট পরিবর্তনেই
হয়। শুধু ফন্টের বদলে লেখার ভাষাও বদলাতে চাইলে Polyglossia-এর ভাষা পরিবর্তনের কমান্ড ব্যবহার করুন।
বন্ধনীসহ সমীকরণের নম্বরের জন্য \texttt{\detokenize{\eqref}}, নামসহ
রেফারেন্সের জন্য \texttt{\detokenize{\autoref}} ব্যবহার করুন। বর্তমান
Polyglossia-এর অঙ্ক রূপান্তরের ত্রুটি সংশোধিত হয়েছে; পুরোনো workaround আর দরকার নেই।

\section*{এক নজরে কমান্ড}
\phantomsection
\addcontentsline{toc}{section}{এক নজরে কমান্ড}
\begin{center}
\begin{tabularx}{\linewidth}{@{}lX@{}}
\toprule
\textbf{কাজ} & \textbf{কমান্ড বা অপশন} \\
\midrule
মূল বাংলা ফন্ট বাছা & \texttt{\detokenize{\setbanglamainfont[features]{font}}} \\
অন্য বাংলা ফন্ট বাছা & \texttt{\detokenize{\setbanglasansfont}}, \texttt{\detokenize{\setbanglamonofont}} \\
ইংরেজি ফন্ট বাছা & \texttt{\detokenize{\setmainfont}}, \texttt{\detokenize{\setsansfont}}, \texttt{\detokenize{\setmonofont}} \\
আগের কাউন্টারের বিন্যাস রাখা & \pkg{changecounternumbering=false} \\
কাউন্টারের মান বাংলায় লেখা & \texttt{\detokenize{\tobangla{counter}}} \\
অঙ্কের সারি রূপান্তর & \texttt{\detokenize{\numtobangla{0012}}} \\
গণিতে বাংলা শব্দ লেখা & \texttt{\detokenize{\text{বাংলা}}} \\
বর্ণভিত্তিক কাউন্টার & \texttt{\detokenize{\alph{counter}}} বা \texttt{\detokenize{\Alph{counter}}} \\
\bottomrule
\end{tabularx}
\end{center}

রূপান্তর চালু থাকলে বর্ণভিত্তিক নম্বর ক থেকে ৎ পর্যন্ত চলে (৩৬টি বর্ণ)।
শূন্যে কিছু লেখা হয় না; ঋণাত্মক বা ৩৬-এর বেশি মানে ``Counter too large'' ত্রুটি আসে।
পরিশিষ্টে পরিশিষ্ট ক-এর মতো নম্বর আসে।

\subsection{নমুনা ও সহায়তা}
সুবিধাগুলো একসঙ্গে দেখতে
\href{https://github.com/Neehan/LaTeXbangla}{রিপোজিটরিতে দেওয়া}
নমুনা ডকুমেন্ট ও ফন্ট নির্দেশিকা
দেখতে পারেন; দুটিরই সোর্স ও PDF দেওয়া আছে। পরীক্ষিত ক্লাসের মধ্যে আছে
article, report, book ও memoir। কোনো কিছু প্রত্যাশামতো কাজ না করলে একটি ছোট
ডকুমেন্ট, ফন্টের নাম, ইঞ্জিনের নাম ও log পাঠান; \texttt{\detokenize{\listfiles}}
প্যাকেজের সংস্করণ লিখে রাখে। আমাদের সঙ্গে যোগাযোগ করতে পারেন
\href{https://github.com/Neehan/LaTeXbangla/issues}{GitHub Issues}-এ বা
\href{https://groups.google.com/g/latexbangla}{\texttt{latexbangla@googlegroups.com}}-এ।

\smallskip
{\small স্বত্ব ২০১৬, ২০২৬ আদীব হাসান।
\href{https://www.latex-project.org/lppl.txt}{LaTeX Project Public License 1.3c বা পরবর্তী সংস্করণে}
প্রকাশিত। এই নির্দেশিকায় বাংলায় Kalpurush ও ইংরেজি মূল লেখায় TeX Gyre Pagella
ব্যবহার করা হয়েছে; প্রথম উদাহরণের ইংরেজি ফন্ট TeX Gyre Adventor।
ফন্টের তুলনায় Noto Sans Bengali ও Siyam Rupali-ও ব্যবহার করা হয়েছে।}
\end{document}
